\section{Observed-law lower bounds}

The estimation and coverage guarantees in \cref{obj:thm:observable-upper,obj:thm:finite-honesty} are complemented by an observed-experiment converse. \Cref{obj:thm:observed-lower} establishes the statistical difficulty of causal evaluation from contaminated-dose records under its stated conditions. Its comparisons concern the sample observed laws \(Q_P^n\) defined in \cref{obj:def:observed-experiment}, and therefore apply directly to procedures based on the available observations. Together with the upper bounds, they supply the lower-bound component of \cref{obj:thm:uniform-frontier}.

The comparison uses the reference and perturbed structural laws \leanref{sym:Palt}{\ensuremath{P_\star,P_+,P_-}} within a bounded Bernoulli witness subclass. These laws share the stratum probabilities, latent-dose design, and Gaussian measurement channel. Their conditional potential means differ through admissible perturbations at the evaluation dose, producing separation of the causal targets. The common design isolates how changes in the conditional mean affect the distribution of observed outcomes and contaminated doses. The structural realization of these witnesses, including the rank coordinate and threshold schedules, is given with the comparison in \cref{obj:thm:observed-lower}.

The statistical argument connects this target separation to closeness of the sample observed laws. When two admissible laws assign separated causal means while generating records that are difficult to distinguish, estimation must incur error on at least one of them. The corresponding interval comparison combines observed-law closeness with the requirement of coverage across the model class. Thus the witness subclass links the causal target to both the estimation criterion in \cref{obj:def:minimax-risk} and the honest-length criterion in \cref{obj:def:honest-length}.

Weighted moment cancellation is the mechanism behind the two inverse-regime comparisons. The packets supplied by \cref{obj:thm:weighted-packet-construction} retain a controlled value at the evaluation point while cancelling moments against the design weight. Gaussian convolution attenuates the resulting signed change in the observed distribution, even as the causal means remain separated. Weighting the cancellation by the latent design makes this mechanism compatible with weak dose support. The two comparisons balance localization and cancellation at the scales relevant to the Gaussian error and the compact dose interval. Their packet formulas, regime-specific choices, and quantitative observed-law bounds are developed in the appendix before use in \cref{obj:thm:observed-lower}.

% Appendix equation labels to assign at their defining displays:
% eq:lower-witness-means -- reference and perturbed conditional means.
% eq:lower-witness-target-separation -- causal-target separation.
% eq:lower-weighted-moment-cancellation -- design-weighted cancellation identity.
% eq:lower-observed-law-comparison -- observed-law divergence comparison.
% eq:lower-gaussian-regime-comparison -- Gaussian inverse-regime tuning and bound.
% eq:lower-compact-regime-comparison -- compact-support inverse-regime tuning and bound.
% Later reader-facing references to these displays must use cleveref.
